% TOP_PROBLEM: {"id": 30006944, "title": "Small Cohen-Macaulay Conjecture", "category_id": 4, "category": {"id": 4, "name": "algebra", "display_name": "Algebra", "description": "Group theory, ring theory, field theory, and algebraic structures.", "slug": "algebra", "order_index": 4, "created_at": "2026-07-31T15:26:25.671Z"}, "status": "open"}
% FIRST_POSTED: 2026-10-06
% ATTEMPT_STATUS: solved
% OPENAI_CLAIM: {"schema_version":1,"source_repo":"https://github.com/openai/math","source_commit":"adc7f1241b42e322a6451854ab7e4b4c146bf78a","release_date":"2026-10-06","match":"full","resolution":"disproved","independently_reviewed":false,"second_pass":"agreed","families":[{"family":"195","title":"A counterexample to the small Cohen–Macaulay module conjecture","subject":"Algebra","manuscripts":[{"title":"A Complete Local Domain without a Small Cohen–Macaulay Module","pdf_path":"preprints/A-Complete-Local-Domain-Without-a-Small-Cohen-Macaulay-Module-September-23-2026/paper.pdf","readme_path":"preprints/A-Complete-Local-Domain-Without-a-Small-Cohen-Macaulay-Module-September-23-2026/README.md","date":"2026-09-23","theorem_ref":"Theorem 1.1"}],"lean":null,"reasoning_trace":null}]}
\documentclass[11pt]{article}
\usepackage[margin=25mm]{geometry}
\usepackage{fontspec}
\setmainfont{Latin Modern Roman}
\usepackage{amsmath,amssymb,mathtools}
\usepackage{unicode-math}
\setmathfont{Latin Modern Math}
\usepackage{enumitem}
\usepackage{xurl,hyperref,bookmark}
\hypersetup{colorlinks=true,urlcolor=blue}
\setcounter{secnumdepth}{0}
\setlength{\parindent}{0pt}
\setlength{\parskip}{5pt}
\setlength{\emergencystretch}{3em}
\begin{document}
\section*{Small Cohen-Macaulay Conjecture}

\subsection{OpenAI's claim}
OpenAI presents a counterexample in the selected manuscript. A normal complete local domain of dimension three over $\mathbb C$ is claimed to admit no nonzero finite module of depth three.

\subsection{Outline of the argument}
The construction completes a polarized surface cone. Extending a hypothetical maximal Cohen--Macaulay module to a regular model, balancing slopes by modifications, and comparing Chern characters yields a numerical surface inequality; an explicit branched surface violates that inequality.

\subsection{Scope relative to this problem}
The proposed complex example is normal and complete, with Krull dimension three. Every nonzero finite module over it has smaller depth, directly opposing universal existence of a small Cohen--Macaulay module. This is presented as a claimed negative answer.

\subsection{Formal verification}
Family 195 has no Lean scope page in this release. Its selected result is not presented as formalized. This site has not executed the Lean code. The scope descriptions do not certify the remaining manuscript arguments.

\subsection{Status caveat}
OpenAI attributes these manuscripts to work by an internal model. Its release cautions that arguments without formalization can contain mistakes. This catalogue has not independently checked the proofs; the labels report the mathematical claims attributed to OpenAI.

\subsection{Sources}
\begin{itemize}
\item OpenAI, \emph{A Complete Local Domain without a Small Cohen–Macaulay Module}, 2026-09-23; citation key \texttt{OAI:A-Complete-Local-Domain-Without-a-Small-Cohen-Macaulay-Module-September-23-2026}. \href{https://github.com/openai/math/blob/adc7f1241b42e322a6451854ab7e4b4c146bf78a/preprints/A-Complete-Local-Domain-Without-a-Small-Cohen-Macaulay-Module-September-23-2026/paper.pdf}{PDF (GitHub)}; \href{https://github.com/openai/math/raw/adc7f1241b42e322a6451854ab7e4b4c146bf78a/preprints/A-Complete-Local-Domain-Without-a-Small-Cohen-Macaulay-Module-September-23-2026/paper.pdf}{PDF (direct)}; \href{https://github.com/openai/math/blob/adc7f1241b42e322a6451854ab7e4b4c146bf78a/preprints/A-Complete-Local-Domain-Without-a-Small-Cohen-Macaulay-Module-September-23-2026/README.md}{Manuscript page and citation}; \href{https://github.com/openai/math/blob/main/preprints/A-Complete-Local-Domain-Without-a-Small-Cohen-Macaulay-Module-September-23-2026/paper.pdf}{Latest version}.
\item \href{https://github.com/openai/math/blob/adc7f1241b42e322a6451854ab7e4b4c146bf78a/CONTENTS.md}{OpenAI catalogue entry}.
\end{itemize}
\end{document}
